\section{Miscellaneous}
\label{appx:misc}

\textbf{Teaser figure description.} We briefly describe how the left hand graph of Figure~\ref{fig:scaling_teaser}, which depicts scaling of task instance collection for the \bugs{} vs. SWE-bench, was created.
For \bugs{}, we simply collected the number of task instances for each repository.
For SWE-bench, we ran the SWE-bench task instance candidate collection script on all $128$ repositories, which first crawls all PRs from a given repository.
Then, each PR that edits at least one or more Python files and changes at least one or more testing related files is converted into a candidate task instance.
Finally, based on the average task instance yield rate reported in~\citet{jimenez_swe-bench_2024}, we estimate the number of viable task instances to be $20$\% of the candidates.
We then determine the number of task instances for \texttt{n} repositories at intervals of $5$ repositories ranging from $5$ to $250$, where the repositories are sorted by number of stars.
In other words, the first five repositories we account for in the figure are the five with the fewest number of stars out of the $128$ repositories used.

\textbf{Extended related works.}
We discuss additional related works briefly, primarily about similar work towards synthesizing trajectories for training LM agents, but for the domain of web tasks.
To improve the interactive capabilities of open source LMs~\citep{chen_fireact_2023}, prior works have also explored trajectory generation techniques for web benchmarks and settings~\citep{xie_osworld_2024,yao_webshop_2023,zhou_webarena_2024}.
For web navigation, existing strategies rely on (1) performing random walks which are then labeled retroactively with instructions~\citep{xiang_language_2023,murty_nnetscape_2024}, (2) using online web tutorials as a source of indirect supervision for generating synthetic trajectories~\citep{ou_synatra_2024}, or (3) collecting human demonstrations~\citep{shen2024scribeagentspecializedwebagents,xu_agenttrek_2024}.
These procedures do not translate well to the software engineering setting; random sequences of command line interactions usually do not achieve meaningful effects on a codebase.
Our cursory efforts around replaying trajectories synthesized from online code edit sequences (e.g. GitHub commit histories) were unsuccessful due to the limited information available, which primarily capture file-level changes without reflecting the underlying skills, decision-making, or the broader context of a software development process.

Our exploration of using SWE-agent to automatically determine installation and testing specifications for a repository is heavily influenced by two research directions - automatic execution environment construction using LMs~\citep{bogin2024superevaluatingagentssetting,eliseeva2025envbenchbenchmarkautomatedenvironment,vergopoulos2025automatedbenchmarkgenerationrepositorylevel}, and generating unit tests using LMs~\citep{mündler2025swtbenchtestingvalidatingrealworld}.
Although relatively much less than SWE-bench style collection, \bugs{} still requires minimal amounts of human labor (around $8$ minutes total per repository).
As we expand \bugs{} to more repositories and languages, we are continuing to consider how to completely automate the environment construction process end to end.
